\scanpage{20}
% 1–4
\deleted{(iii) Etude du cas où $\deleted{\unsure{F=|F\text{ est}}}\left\{\begin{aligned}a_i&=0\quad\text{pour }i\geq N+1\\ a_N&=1\end{aligned}\right.$}

\deleted{i.e. du cas où}
\[
\deleted{F=Z^N+a_{N-1}Z^{N\unsure{-|+}1}\ldots+a_0}
\]
\deleted{les $a_i$ $(0\leq i\leq N-1)$ nilpotents.}

% 5–12
\[
\deleted{\underset{\substack{\text{base }(Z^i)_{0\leq i\leq N-1}\\\text{sur }A}}{A[Z]/FA[Z]}\longrightarrow A[[Z]]/FA[[Z]]}
\]
\deleted{est alors surjectif (car les $Z^i$ $(0\leq i\leq N-1)$ engendrent $A[[Z]]/FA[[Z]]$ comme $A$-module) pr dire que c'est bijectif i.e. que $(FA[[Z]])\cap A[Z]=FA[Z]$ (il revient au $=$ de dire que les \deleted{\unsure{deux classes|classes de}} $Z^i$ $(0\leq i\leq N-1)$ forment une base de $A[[Z]]/FA[[Z]]$).}

% 13–14
\deleted{Soit donc $a_0+a_1Z\ldots+a_{N-1}Z^{N-1}$ qui soit multiple de $F$ dans $A[[Z]]$, \unsure{soi.|soit|i.e.}}

% 15–17
Considérons l'application
\[
\begin{aligned}
\varphi:\quad A^N+A[[Z]]&\longrightarrow A[[Z]]\\
\big((b_i)_{0\leq i\leq N-1},Q\big)&\longrightarrow\sum_{0\leq i\leq N-1}\deleted{\unsure{Q_i|Q}}b_iZ^i+QF
\end{aligned}
\]

% 18–23
l'application déduite par passage aux gradués associés pour la filtration $\mathfrak m$-adique (finie) est
\[
(\mathrm{Gr}\,A)^N\times(\mathrm{Gr}\,A)[[Z]]\longrightarrow(\underline{\unsure{\mathrm{Gr}'|\mathrm{Gr}}\,A})[[Z]]
\]
\unsure{donné|donnée} \deleted{par} en degré $n$ par
\[
((b_i),Q)\longrightarrow\sum_{0\leq i\leq N-1}b_iZ^i+Q\dot F
\]

% 24–26
et comme $\dot F=\deleted{\dot a_{N+1}}\alpha Z^N$ $(\alpha\unsure{=|-}a_{N+1}\in A_0^*)$ ceci est manifestement bijectif, donc $\varphi$ aussi, cqfd.
% archive sheet number: 19
